% !TeX root = ../main.tex
\section{Koncepcja bistatycznego radaru pasywnego}
\label{sec:pbr}

% ------------------------------------------------------------------------
% Geometria bistatyczna:
%   - baza bistatyczna L, opóźnienie bistatyczne, przesunięcie częstotliwości Dopplera
% Zjawiska fizyczne w pasywnej radiolokacji:
%   - przesłuch bezpośredni (DPI), echa celów ruchomych
%   - stacjonarny clutter wielodrogowy: kanał syntetyczny vs model wewnątrzbudynkowy TGax Model-B
% Etapy przetwarzania sygnału:
%   - akwizycja, tłumienie składowej stałej/DPI, kompresja odległość-Doppler, detekcja celów
% Metody estymacji: funkcja niejednoznaczności (CAF) a odpowiedź kanału (CFR/CIR):
%   - zalety paradygmatu estymacji kanału (Zero-Forcing) w systemach OFDM
% Zalety i ograniczenia pasywnego radaru Wi-Fi:
%   - niska moc nadawcza AP, kompromisy rozdzielczości i prędkości jednoznacznej
% ------------------------------------------------------------------------

\subsection{Geometria bistatyczna radaru pasywnego}
\label{subsec:pbr_geometry}

\subsection{Zjawiska fizyczne w pasywnej radiolokacji}
\label{subsec:pbr_physics}

\subsection{Etapy przetwarzania sygnału w radarze pasywnym}
\label{subsec:pbr_processing_stages}

\subsection{Metody estymacji: funkcja niejednoznaczności (CAF) a odpowiedź kanału (CFR/CIR)}
\label{subsec:pbr_caf_vs_cir}

\subsection{Zalety i ograniczenia pasywnego radaru Wi-Fi}
\label{subsec:pbr_pros_cons}
